% 1. introduzione
\section{Resoconto}

\subsection{Scopo del documento}

Lo scopo del documento è descrivere l'analisi dei requisiti 
del progetto \textbf{Second Brain}. Il testo definisce in modo
strutturato i casi d’uso e le funzionalità richieste,
classificandole in requisiti obbligatori, desiderabili
e opzionali, secondo le specifiche fornite da Zucchetti. 


\subsection{Riferimenti}

\begin{itemize}
	\subsubsection{Riferimenti normativi}
        \item Capitolato C6 - Progetto Second Brain: \\  \url{https://www.math.unipd.it/~tullio/IS-1/2025/Progetto/C6.pdf}
        \item Regolamento progetto didattico: \\ \url{https://www.math.unipd.it/~tullio/IS-1/2025/Dispense/PD1.pdf}
        \item \href{https://github.com/Synergo-Unit-UniPD/Documentary-Repo/blob/main/docs/CC/doc_gestionali/doc_interni/norme_di_progetto/norme_di_progetto.pdf}{Norme di Progetto v1.1.2}
    \end{itemize}
    \subsubsection{Riferimenti informativi}
    \begin{itemize}
        \item \href{https://synergo-unit-unipd.github.io/Documentary-Repo/}{Verbali}
        \item \href{https://github.com/Synergo-Unit-UniPD/Documentary-Repo/blob/main/docs/RTB/doc_gestionali/doc_interni/glossario/glossario.pdf}{Glossario v2.1.0}
\end{itemize}

\newpage

\section{Descrizione del progetto}

\subsection{Obbiettivi del progetto}
Il prodotto finale prevede lo sviluppo di un’applicazione
in grado di gestire la redazione di testi con sintassi
Markdown, garantirne la corretta resa formattata e
integrare funzionalità basate su modelli linguistici
di grandi dimensioni. In particolare, Second Brain
utilizzerà LLM per offrire varie operazioni tra
cui di sintesi, riformulazione e generazione di contenuti, 
oltre a un modulo di valutazione ispirato al metodo dei 
“Sei Cappelli per Pensare” di Edward De Bono.

\subsection{Funzionalità del progetto}

\begin{itemize}
    \item \textbf{Editing di testo con marcatori MarkDown}:
	L'utente potrà scrivere in un'area che accetta
    testo e marcatori in formato MarkDown;
    \item \textbf{Utilizzo LLM}; L'utente, 
	tramite supporto di un LLM, potrà riassumere,
	riscrivere o tradurre il testo;
    \item \textbf{Creazione intero documento tramite LLM}: 
	L'utente potrà far sviluppare l'intero testo all'LLM stesso
	sulla base di del modo di pensare scelto,
	ispirato ai “Sei Cappelli per Pensare” di Edward
	De Bono.;
\end{itemize}

\newpage

\section{Casi d'uso}


